%% kapa: exam=1 round=21 year=2010 period=2 form=A subject=accounting subject_ko=회계학 seq=12 number=12
%% source: Q-Net 감정평가사 기출문제 (cst00309) artlSeq=5024181 file=extracted/1008/5024181/2010년 제21회 감정평가사 1차시험 기출문제/2교시.hwp
%% conversion: kordoc 4.24.51 -> tools/gen1.py (auto); review=proofread-gpt-sol; flags=none
\begin{question}{12}
주당이익 계산에 있어서 당해 기간 및 비교 표시되는 모든 기간의 가중평균유통보통주식수는 잠재적보통주의 전환을 제외하고, 상응하는 자원의 변동 없이 유통보통주식수를 변동시키는 사건을 반영하여 조정한다. 다음 중 상응하는 자원의 변동 없이 유통보통주식수를 변동시키는 사례가 아닌 것은?
\begin{choices}
  \item 다른 거래 없이 1주당 액면 ￦100인 주식 5주를 액면 ￦500인 주식 1주로 병합하였다.
  \item 다른 거래 없이 1주당 액면 ￦5,000인 주식 1주를 액면 ￦500인 주식 10주로 분할하였다.
  \item 이익준비금을 자본금으로 전입하였다.
  \item 주식발행초과금을 자본금으로 전입하였다.
  \item 기존 주주로부터의 차입금을 자본으로 전환하였다.
\end{choices}
\end{question}
\answer{5}
